% MG06. En exc:paley-codigo, tras definir C y chi, antes de la frase
% «Las identidades de la suma cuadrática...» y de reducir módulo tres.
L'identité quadratique de cette coordinatisation est obtenue par une
corrélation finie. On conserve l'ordre
\((\infty,0,1,2,3,4)\), ainsi que l'identification des six
positions et le caractère \(\chi\) déjà déclarés.

\begin{lemma}[Corrélation du caractère quadratique en cinq positions]
\label{mg06:correlacion}
Pour \(a,b\in\mathbb F_5\) avec \(a\ne b\), on a
\begin{equation}
 \sum_{t\in\mathbb F_5}\chi(t-a)\chi(t-b)=-1.
 \label{mg06:identidad-correlacion}
\end{equation}
\end{lemma}
\begin{proof}
La substitution \(u=(t-a)/(b-a)\) est bijective, puisque \(b-a\) est
inversible dans \(\mathbb F_5\). De plus,
\[
 (t-a)(t-b)=(b-a)^2u(u-1),\qquad
 \chi((b-a)^2)=1.
\]
La multiplicativité de \(\chi\), y compris son prolongement par zéro, réduit la somme à \(\sum_u\chi(u(u-1))\). Les cinq évaluations sont
\begin{equation}
\begin{array}{c|rrrrr}
 u&0&1&2&3&4\\\hline
 u(u-1)\pmod5&0&0&2&1&2\\
 \chi(u(u-1))&0&0&-1&1&-1
\end{array}
\label{mg06:cinco-evaluaciones}
\end{equation}
Leur somme est \(-1\), ce qui démontre
\eqref{mg06:identidad-correlacion}.
\end{proof}

\begin{proposition}[Orthogonalité de la coordinatisation à six positions]
\label{mg06:conferencia}
La matrice \(C\) satisfait
\begin{equation}
 C=C^{\mathsf T},\qquad CC^{\mathsf T}=5I_6,
 \qquad C^2=5I_6.
 \label{mg06:identidad-conferencia}
\end{equation}
\end{proposition}
\begin{proof}
Puisque \(\chi(-1)=\chi(4)=1\), on a \(\chi(b-a)=\chi(a-b)\) ; les entrées faisant intervenir \(\infty\) sont également symétriques. Chaque ligne contient une entrée nulle et cinq entrées de module un, de sorte que le carré de sa norme est cinq. Pour deux lignes finies distinctes, le produit scalaire reçoit \(+1\) de la coordonnée \(\infty\) et \(-1\) des cinq coordonnées finies par \eqref{mg06:identidad-correlacion} ; leur somme est zéro. Le produit scalaire de la ligne infinie avec une ligne finie est \(\sum_{t\in\mathbb F_5}\chi(t)=0+1-1-1+1=0\). Toutes les entrées de \(CC^{\mathsf T}\) sont ainsi évaluées. La symétrie transforme cette identité en \(C^2=5I_6\).
\end{proof}

L'argument conserve l'ordre de la coordinatisation qui sera réduite modulo
trois. L'identité quadratique et le code graphique ultérieur
utilisent donc les mêmes coordonnées, avec leurs incidences et leurs supports.
